% ============================================================
%  BACKWARD TRAJECTORY COMPLETION IN BOUNDED PHASE SPACE:
%  S-ENTROPY FRAMEWORK, VIRTUAL SUB-STATES, AND THE LHC
%  TRIGGER AS A CATEGORICAL NAVIGATOR
% ============================================================

\documentclass[twocolumn,10pt,a4paper]{article}

\usepackage[utf8]{inputenc}
\usepackage[T1]{fontenc}
\usepackage{lmodern}
\usepackage[a4paper,top=2cm,bottom=2.2cm,left=1.8cm,right=1.8cm,columnsep=0.6cm]{geometry}
\usepackage{amsmath,amssymb,amsthm,mathtools}
\usepackage{bm}
\usepackage{booktabs}
\usepackage{array}
\usepackage{enumitem}
\usepackage{graphicx}
\usepackage{xcolor}
\usepackage{microtype}
\usepackage{mathrsfs}
\usepackage{algorithm}
\usepackage{algorithmic}
\usepackage[round,sort&compress]{natbib}
\usepackage[colorlinks=true,linkcolor=blue!60!black,
            citecolor=blue!60!black,urlcolor=blue!60!black]{hyperref}
\usepackage{fancyhdr}

\pagestyle{fancy}
\fancyhf{}
\fancyhead[LE,RO]{\thepage}
\fancyhead[RE]{\small Backward Trajectory Completion in Bounded Phase Space}
\fancyhead[LO]{\small Sachikonye}
\renewcommand{\headrulewidth}{0.4pt}

% ---- Theorem environments -------------------------------------------
\newtheoremstyle{thmstyle}{\topsep}{\topsep}{\itshape}{}{\bfseries}{.}{.5em}{}
\theoremstyle{thmstyle}
\newtheorem{theorem}{Theorem}[section]
\newtheorem{lemma}[theorem]{Lemma}
\newtheorem{proposition}[theorem]{Proposition}
\newtheorem{corollary}[theorem]{Corollary}
\theoremstyle{definition}
\newtheorem{definition}[theorem]{Definition}
\newtheorem{axiom}{Axiom}
\newtheorem{example}[theorem]{Example}
\theoremstyle{remark}
\newtheorem{remark}[theorem]{Remark}

% ---- Custom commands ------------------------------------------------
\newcommand{\R}{\mathbb{R}}
\newcommand{\N}{\mathbb{N}}
\newcommand{\Zp}{\mathbb{Z}^{+}}
\newcommand{\tP}{t_{\mathrm{P}}}
\newcommand{\kB}{k_{\mathrm{B}}}
\newcommand{\Parts}{\mathcal{P}}
\newcommand{\Sfl}{S_{\flat}}
\newcommand{\Sk}{S_{k}}
\newcommand{\St}{S_{t}}
\newcommand{\Se}{S_{e}}
\newcommand{\Svec}{\mathbf{S}}
\newcommand{\Sspace}{\mathcal{S}}
\newcommand{\receiver}{\mathcal{R}}
\newcommand{\Osc}{\mathbf{Osc}}
\newcommand{\Cat}{\mathbf{Cat}}
\newcommand{\Part}{\mathbf{Part}}
\newcommand{\catpow}{\kappa}
\newcommand{\nP}{n_{\mathrm{P}}}
\newcommand{\etaC}{\eta_{C}}

% =========================================================
\title{\textbf{Backward Trajectory Completion in Bounded Phase Space:\\
S-Entropy Framework, Virtual Sub-States, and\\
the LHC Trigger as a Categorical Navigator}}

\author{
Kundai Farai Sachikonye\\
Technical University of Munich\\
Chair of Proteomics and Bioanalytics\\
\texttt{kundai.sachikonye@wzw.tum.de}
}
\date{\today}
% =========================================================

\begin{document}
\maketitle

% =========================================================
\begin{abstract}
\noindent
The LHC trigger faces a classification task: given a pattern of detector
hits (the \emph{endpoint}), determine the originating physics process (the
\emph{root}). We prove that this task is precisely backward trajectory
completion in a ternary refinement hierarchy, and that it can be solved in
$\mathcal{O}(\log_3 N)$ steps via the S-entropy embedding into $[0,1]^3$
with the product Fisher metric — compared to $\Theta(N)$ for any algorithm
restricted to physically realizable intermediate states.

The key mechanism is \emph{virtual sub-states}: intermediate decompositions
$(s_1, s_2, s_3) \in \R^3$ with mean $\bar s = s \in [0,1]$ but at least
one component outside $[0,1]^3$. We prove three results about virtual
sub-states that are central to trigger design:
(i) they exist in positive Lebesgue measure (Theorem~\ref{thm:virtual_existence});
(ii) restricting to physically realizable sub-states collapses backward
completion to $\Theta(N)$ (Theorem~\ref{thm:collapse});
(iii) two backward trajectories sharing an endpoint are observationally
indistinguishable from any metric invariant at the endpoint (Path Opacity,
Theorem~\ref{thm:path_opacity}).

We establish the S-entropy framework from which these results derive: a
bounded receiver $\receiver = (\Sigma, \Phi, K, \mathrm{dec})$ has a
strictly positive alignment floor $\Sfl(\receiver) > 0$
(Theorem~\ref{thm:floor}), and three algebraic structures
(oscillatory $\Osc$, categorical $\Cat$, partition $\Part$) are mutually
equivalent via explicit conversion functors
(Theorem~\ref{thm:triple_equiv}).

The paper concludes with the Aperture Identification theorem
(Theorem~\ref{thm:aperture_id}): geometric aperture, virtual sub-state,
miracle subtask, information catalyst, and Maxwell demon are five names
for the same mathematical object — the partition-theoretic mechanism
that enables $\mathcal{O}(\log N)$ navigation through the state space
of physics events. This identifies the LHC trigger as a physical
implementation of the categorical aperture, operating on collision events
rather than molecular spectra.

\medskip
\noindent\textbf{Keywords:} backward trajectory completion; S-entropy;
virtual sub-states; ternary refinement hierarchy; LHC trigger; triple
equivalence; path opacity; categorical aperture; information catalyst;
partition algebra
\end{abstract}

\tableofcontents

% =========================================================
\section{Introduction}
\label{sec:intro}
% =========================================================

\subsection{The classification problem}

The LHC trigger must classify $4 \times 10^7$ collision events per second.
Each event is characterised by a pattern of detector hits: energies
deposited in calorimeter cells, track momenta in the inner detector, muon
hits in the outer spectrometer. The trigger must map this hit pattern to
a class label (``interesting physics'' or ``background'') within 2.5~$\mu$s.

Framed abstractly: the trigger is given the \emph{endpoint} of a trajectory
through the partition state space of the detector (the hit pattern) and must
recover the \emph{root} (the originating physics process). This is backward
trajectory completion.

The conventional approach — forward simulation and threshold comparison — is
computationally equivalent to forward trajectory construction: build all
possible event topologies from assumed physics models and compare with the
observed hit pattern. This is $\Theta(N)$ where $N$ is the number of
possible physics processes, which is astronomically large.

We show that backward trajectory completion, using the S-entropy embedding
and virtual sub-states, achieves $\mathcal{O}(\log_3 N)$.

\subsection{Outline and main results}

Section~\ref{sec:receiver} develops the S-entropy framework: bounded
receivers, the S-functional, and the Floor Positivity theorem.
Section~\ref{sec:triple} proves the Triple Equivalence between oscillatory,
categorical, and partition descriptions.
Section~\ref{sec:ternary} defines the ternary refinement hierarchy and
the S-entropy embedding.
Section~\ref{sec:backward} proves backward navigation complexity and
establishes the collapse theorem.
Section~\ref{sec:virtual} develops the theory of virtual sub-states.
Section~\ref{sec:opacity} proves path opacity and its trigger implications.
Section~\ref{sec:trigger} applies the theory to the LHC trigger: the
detector as a backward navigator, the cell registry as an aperture,
and the working example of the di-muon trigger.
Section~\ref{sec:aperture} proves the Aperture Identification theorem,
unifying all five names for the same mechanism.

% =========================================================
\section{S-Entropy Framework and Receiver Theory}
\label{sec:receiver}
% =========================================================

\subsection{Bounded receivers}

\begin{definition}[Candidate and Truth Spaces]
A \emph{candidate space} $\mathcal{X}$ is a non-empty set of admissible
solutions. A \emph{truth space} $\mathcal{X}^* \subseteq \mathcal{X}$
is the designated set of correct solutions. We allow $|\mathcal{X}^*| > 1$
to accommodate operationally indistinguishable solutions.
\end{definition}

\begin{definition}[Bounded Receiver]
\label{def:receiver}
A \emph{receiver} over $(\mathcal{X}, \mathcal{X}^*)$ is a quadruple
$\receiver = (\Sigma, \Phi, K, \mathrm{dec})$ where:
\begin{enumerate}[nosep,label=(\roman*)]
\item $\Sigma$ is the \emph{signal space};
\item $\Phi : \Sigma \to K$ is the \emph{decoder};
\item $K$ is the \emph{knowledge framework} with operations $\{\circ_j\}$;
\item $\mathrm{dec} : K \to \mathcal{X}$ is the \emph{candidate projection}.
\end{enumerate}
The receiver is \emph{bounded} if $|K| < |\mathcal{X}|$ (strictly smaller
than the candidate space).
\end{definition}

\begin{definition}[S-Functional]
\label{def:s_functional}
An \emph{S-functional} is a map
$\mathcal{S} : \mathfrak{R} \times \mathcal{X} \times \mathcal{X}^*
\to [0, 100]$
measuring misalignment: $\mathcal{S}_\receiver(x, x^*) = 0$ denotes perfect
alignment; $\mathcal{S}_\receiver(x, x^*) = 100$ denotes maximum misalignment.
It satisfies four axioms:
\begin{enumerate}[nosep,label=(S\arabic*)]
\item Non-negativity: $\mathcal{S}_\receiver(x,x^*) \geq 0$.
\item Boundedness: $\mathcal{S}_\receiver(x,x^*) \leq 100$.
\item Self-truth lower bound: $\inf_x \mathcal{S}_\receiver(x,x^*) = \Sfl(\receiver)$
for some function $\Sfl : \mathfrak{R} \to (0,100]$.
\item Receiver coherence: $\mathcal{S}_\receiver(x,x^*)
= \delta_\receiver(\Phi(\sigma_x), \Phi(\sigma_{x^*}))$
for a non-decreasing $\delta_\receiver$ with $\delta_\receiver(\sigma,\sigma) = \Sfl(\receiver)$.
\end{enumerate}
\end{definition}

\subsection{The Floor Positivity theorem}

\begin{theorem}[Floor Positivity]
\label{thm:floor}
For every bounded receiver $\receiver$, $\Sfl(\receiver) > 0$.
\end{theorem}

\begin{proof}
Suppose $\Sfl(\receiver) = 0$. By (S3), there exists $x_0 \in \mathcal{X}$
with $\mathcal{S}_\receiver(x_0, x^*) = 0$. By (S4), this requires
$\delta_\receiver(\Phi(\sigma_{x_0}), \Phi(\sigma_{x^*})) = 0$.
Since $\delta_\receiver(\sigma,\sigma) = \Sfl = 0$ and $\delta_\receiver$ is
non-decreasing, $\Phi(\sigma_{x_0}) = \Phi(\sigma_{x^*})$ in $K$.
The receiver is bounded: $|K| < |\mathcal{X}|$, so
$|\mathrm{dec}^{-1}(\{x^*\})| \leq |K| < |\mathcal{X}|$, leaving
candidates unreachable in $K$. This contradicts the requirement that every
candidate has a representative in $K$. Hence $\Sfl(\receiver) > 0$.
\end{proof}

\begin{corollary}[Smallest Knowable Error]
\label{cor:floor}
No candidate achieves $\mathcal{S}_\receiver(x, x^*) < \Sfl(\receiver)$.
The floor is intrinsic to the receiver: it depends on $|K|$, $\delta_\receiver$,
and $\mathrm{dec}$, not on the candidate or the truth.
\end{corollary}

\subsection{S-entropy coordinates}

\begin{definition}[S-Entropy Coordinate Space]
\label{def:sentropy}
The \emph{S-entropy coordinate space} is the unit cube
$\Sspace = [0,1]^3$ with coordinates $(\Sk, \St, \Se)$ representing:
\begin{itemize}[nosep]
\item $\Sk$ (\emph{knowledge entropy}): normalised Shannon entropy of the
frequency distribution of modes — how uniformly the knowledge is spread.
\item $\St$ (\emph{temporal entropy}): relative entropy of the phase position
within the oscillation cycle.
\item $\Se$ (\emph{evolution entropy}): entropy of the trajectory history —
which sequence of states has been visited.
\end{itemize}
All three coordinates lie in $[0,1]$ by normalisation; the cube $\Sspace$
is a compact metric space under the Euclidean metric.
\end{definition}

\begin{theorem}[Ternary Encoding of S-Space]
\label{thm:ternary_encoding}
Every point $\Svec \in \Sspace$ corresponds to a unique infinite ternary
string $\{t_i\}_{i \geq 1}$ with $t_i \in \{0,1,2\}$ encoding:
$t_i = 0$ refines $\Sk$; $t_i = 1$ refines $\St$; $t_i = 2$ refines $\Se$.
A $k$-trit string addresses one cell in the $3^k$-fold partition of $\Sspace$.
As $k \to \infty$, the addressed cell converges to a unique point.
\end{theorem}

\begin{proof}
At each refinement step, the trit $t_i$ selects which of three equal
sub-intervals along the indicated axis contains the point.
After $k$ steps with $n_j$ trits assigned to axis $j$ ($j \in \{0,1,2\}$),
the $j$-th coordinate is pinned to an interval of width $3^{-n_j}$.
Since $n_j \to \infty$ as $k \to \infty$ (each axis is refined at least
$k/3$ times by the recurrence structure of bounded oscillations),
all three interval widths $\to 0$ and the intersection is a single point.
\end{proof}

% =========================================================
\section{Triple Equivalence}
\label{sec:triple}
% =========================================================

\subsection{Three algebraic structures}

\begin{definition}[Oscillatory Structure]
An \emph{oscillatory structure} $\mathfrak{O} = (M, \omega, \Phi_\mathfrak{O}, \mu)$
consists of: mode set $M$; frequency assignment $\omega : M \to \R_{\geq 0}$;
phase function $\Phi_\mathfrak{O}(m,t) : M \times \R \to S^1$ with
$\Phi_\mathfrak{O}(m, t+s) = \Phi_\mathfrak{O}(m,t) \cdot e^{2\pi i \omega(m) s}$;
probability measure $\mu$ on $M$.
Morphisms preserve frequencies and phase compatibility.
The category is $\Osc$.
\end{definition}

\begin{definition}[Categorical Structure]
A \emph{categorical structure} $\mathfrak{C} = (C, \prec, \lambda, \nu)$
consists of: cell set $C$; strict partial order $\prec$ (refinement);
labelling $\lambda : C \to \Lambda$; measure $\nu$.
Morphisms preserve $\prec$ and $\lambda$. The category is $\Cat$.
\end{definition}

\begin{definition}[Partition Structure]
A \emph{partition structure} $\mathfrak{P} = (\Omega, \mathscr{P}, \pi)$
consists of: measurable space $\Omega$; refining sequence
$\mathscr{P} = \{P_n\}_{n \in \N}$ of finite measurable partitions;
measure $\pi$ with $\pi(\partial p) = 0$ for each $p \in P_n$.
Morphisms preserve refinement and measure. The category is $\Part$.
\end{definition}

\subsection{Conversion functors}

\begin{definition}[Conversion Functors]
\label{def:functors}
Define:
\begin{enumerate}[nosep,label=(\roman*)]
\item $F_{OC} : \Osc \to \Cat$: sends $(M, \omega, \Phi_\mathfrak{O}, \mu)$
to $(M, \prec_\omega, \lambda_\omega, \mu)$ where $m_1 \prec_\omega m_2$ iff
$\omega(m_1) | \omega(m_2)$ (rational ratio with bounded denominator), and
$\lambda_\omega(m) = (\lfloor\log_2\omega(m)\rfloor, \arg\Phi_\mathfrak{O}(m,0))$.
\item $F_{CP} : \Cat \to \Part$: sends $(C, \prec, \lambda, \nu)$ to
$(C/{\sim}, \{P_n\}, \pi)$ where $\sim_n$ identifies cells agreeing on
the first $n$ label coordinates.
\item $F_{PO} : \Part \to \Osc$: sends $(\Omega, \mathscr{P}, \pi)$ to
$(P_\infty, \omega_\mathfrak{P}, \Phi_\mathfrak{P}, \pi)$ with
$\omega_\mathfrak{P}(p) = 1/\pi(p)$ (Kac's lemma analogue).
\end{enumerate}
\end{definition}

\begin{theorem}[Triple Equivalence]
\label{thm:triple_equiv}
The functors $F_{OC}$, $F_{CP}$, $F_{PO}$ are part of mutually inverse
equivalences of categories:
\begin{align}
F_{CO} \circ F_{OC} &\cong \mathrm{id}_\Osc, &
F_{OC} \circ F_{CO} &\cong \mathrm{id}_\Cat, \\
F_{PC} \circ F_{CP} &\cong \mathrm{id}_\Cat, &
F_{CP} \circ F_{PC} &\cong \mathrm{id}_\Part, \\
F_{OP} \circ F_{PO} &\cong \mathrm{id}_\Part, &
F_{PO} \circ F_{OP} &\cong \mathrm{id}_\Osc.
\end{align}
Composing two of the three yields the third up to natural isomorphism.
\end{theorem}

\begin{proof}[Proof sketch]
The inverse $F_{CO}$ is constructed from $F_{OC}$ by recovering the
frequency from the first label coordinate ($\omega(c) = 2^{\lambda_1(c)}$)
and the phase from the second ($\Phi(c,0) = e^{i\lambda_2(c)}$).
The composition $F_{CO} \circ F_{OC}$ acts on $(M, \omega, \Phi_\mathfrak{O}, \mu)$
and recovers the structure up to the discretisation error of $\lfloor\cdot\rfloor$,
which vanishes as the labelling refines. Naturality is verified pointwise;
the triangle identities follow from standard category theory
\citep{maclane1998}.
The cyclic identities $F_{OC} \cong F_{PC} \circ F_{OP}$ etc. follow from
uniqueness of equivalences up to natural isomorphism: any two functors
inducing the same equivalence are naturally isomorphic.
\end{proof}

\begin{corollary}[Optimal Representation]
\label{cor:optimal_rep}
For any computation on a bounded oscillatory system, there exist equivalent
formulations in all three representations $(\Osc, \Cat, \Part)$.
The computation may be performed in whichever representation makes it
computationally cheapest.
\end{corollary}

\begin{remark}[Physical meaning]
Triple Equivalence is not merely algebraic. It says that measuring
any one of: the oscillation frequency, the categorical state label, or the
partition cell address — fully determines the other two. For the LHC trigger,
this means: measuring the detector hit pattern (partition description)
uniquely determines the oscillatory state (particle kinematics) and the
categorical state (event class). The three representations are three windows
onto the same physical reality.
\end{remark}

% =========================================================
\section{Ternary Refinement Hierarchy and S-Entropy Embedding}
\label{sec:ternary}
% =========================================================

\subsection{The hierarchy}

\begin{definition}[Ternary Refinement Hierarchy]
\label{def:ternary_hierarchy}
A \emph{ternary refinement hierarchy} of depth $k$ is a partition structure
$\mathfrak{P} = (\Omega, \{P_j\}_{j=0}^k, \pi)$ where:
$|P_0| = 1$ (the trivial partition),
$|P_{j+1}| = 3|P_j|$ (each cell splits into exactly 3 children),
and each block of $P_j$ is a disjoint union of 3 blocks of $P_{j+1}$.
The leaves are $P_k$ with $|P_k| = 3^k = N$ elements.
\end{definition}

\begin{definition}[Trajectory in the Hierarchy]
\label{def:trajectory}
A \emph{trajectory} in a depth-$k$ hierarchy is a path
$\mathcal{T} = (p_0, p_1, \ldots, p_k)$ with $p_0 \in P_0$,
$p_{j+1}$ a child of $p_j$ in $P_{j+1}$.
The \emph{endpoint} is $\partial(\mathcal{T}) = p_k \in P_k$.
The \emph{root} is $p_0$.
\end{definition}

\begin{definition}[Completion Problem]
\label{def:completion}
The \emph{trajectory completion problem} takes as input the hierarchy
$\mathfrak{P}$ and an endpoint $p_k \in P_k$, and recovers the unique
trajectory $\mathcal{T}$ with $\partial(\mathcal{T}) = p_k$.
\end{definition}

\subsection{The S-entropy embedding}

\begin{definition}[S-Entropy Embedding]
\label{def:s_embed}
The \emph{S-entropy embedding} of a depth-$k$ hierarchy is
$\iota_k : P_k \to [0,1]^3$ sending each leaf $p \in P_k$ to
$\Svec(p) = (\Sk(p), \St(p), \Se(p))$ given by the base-3 expansion:
\[
\Sk(p) = \sum_{j=1}^k \frac{a_j^k(p)}{3^j}, \quad
\St(p) = \sum_{j=1}^k \frac{a_j^t(p)}{3^j}, \quad
\Se(p) = \sum_{j=1}^k \frac{a_j^e(p)}{3^j},
\]
where $a_j^c(p) \in \{0,1,2\}$ is the address digit of $p$ at level $j$
in the $c$-th coordinate.
\end{definition}

\begin{definition}[Product Fisher Metric]
The \emph{product Fisher metric} on $[0,1]^3$ is the Riemannian metric
\[
g(\Svec) = \sum_{c \in \{k,t,e\}} \frac{dS_c^2}{S_c(1-S_c)}.
\]
\end{definition}

\begin{theorem}[Embedding Necessity for Sub-linear Completion]
\label{thm:embedding_necessity}
Any algorithm solving trajectory completion in $o(N)$ queries requires
a metric on the leaf set $P_k$ that distinguishes leaves at distance
proportional to $\log_3|P_k|$. The S-entropy embedding with the product
Fisher metric provides exactly this.
\end{theorem}

\begin{proof}
By the information-theoretic lower bound, any comparison-based
algorithm on $N = 3^k$ leaves without geometric structure requires
$\Omega(N)$ comparisons (each query reveals one bit of structure;
there are $\Omega(N \log_2 3)$ bits to discover). With a metric whose
ball of radius $r$ contains at most $f(r)$ leaves, navigation requires
$\mathcal{O}(\log f^{-1}(N))$ steps. For the S-entropy embedding with
Fisher metric, the ball of radius $r$ contains at most $\mathcal{O}(r^3 3^k)$
leaves (three-dimensional ball volume times leaf density), yielding
$\mathcal{O}(\log_3 N)$ steps.
\end{proof}

\begin{theorem}[Recursive Self-Similarity]
\label{thm:self_similar}
The S-entropy embedding is strictly self-similar: for any block $b \in P_j$,
the embedding restricted to leaves under $b$ is isometric (after rescaling
by $3^{-(k-j)}$) to the embedding of a depth-$(k-j)$ hierarchy.
\end{theorem}

\begin{proof}
Leaves under $b$ share the first $j$ address digits; the remaining
$k-j$ digits enumerate the sub-tree below $b$. Truncating the shared
prefix and rescaling by $3^{-(k-j)}$ recovers the full depth-$(k-j)$ structure
with the same Fisher metric.
\end{proof}

% =========================================================
\section{Backward Navigation Complexity}
\label{sec:backward}
% =========================================================

\subsection{Penultimate state and backward steps}

\begin{definition}[Penultimate State]
\label{def:penultimate}
For a trajectory $\mathcal{T} = (p_0, p_1, \ldots, p_k)$ with $k \geq 1$,
the \emph{penultimate state} is $p_{k-1}$ — the unique parent of
$p_k$ in $P_{k-1}$.
\end{definition}

\begin{theorem}[Penultimate Existence and Uniqueness]
\label{thm:penultimate}
For every leaf $p_k \in P_k$ with $k \geq 1$, the penultimate state
$p_{k-1}$ exists and is unique.
\end{theorem}

\begin{proof}
Existence: the hierarchy is a rooted tree; every non-root node has a parent.
Uniqueness: each block belongs to exactly one parent block (partitions are
disjoint), so the parent is unique.
\end{proof}

\begin{theorem}[Backward Navigation Complexity]
\label{thm:backward_complexity}
The backward navigation algorithm from any leaf $p_k$ to the root $p_0$
visits exactly $\log_3 N = k$ states and runs in $\mathcal{O}(\log_3 N)$ steps.
\end{theorem}

\begin{proof}
Each backward step from $p_j$ to $p_{j-1}$ is implemented by dropping
the last trit of the base-3 expansion of $\iota_j(p_j)$ — an $\mathcal{O}(1)$
operation. Starting at $p_k$ and ending at $p_0$ requires exactly $k = \log_3 N$
steps.
\end{proof}

The backward navigation algorithm is:

\begin{algorithm}[h]
\caption{Backward trajectory completion}
\begin{algorithmic}[1]
\REQUIRE Hierarchy $\mathfrak{P}$, leaf $p_k$
\ENSURE Full trajectory $(p_0, p_1, \ldots, p_k)$
\STATE $\mathcal{T} \leftarrow [p_k]$
\STATE $j \leftarrow k$
\WHILE{$j > 0$}
  \STATE $p_{j-1} \leftarrow \textsc{Parent}(\mathfrak{P}, p_j)$
    \COMMENT{$\mathcal{O}(1)$ via base-3 truncation}
  \STATE $\mathcal{T}.\textsc{prepend}(p_{j-1})$
  \STATE $j \leftarrow j - 1$
\ENDWHILE
\RETURN $\mathcal{T}$
\end{algorithmic}
\end{algorithm}

\subsection{Forward asymmetry}

\begin{theorem}[Forward Asymmetry]
\label{thm:forward_asymmetry}
Forward trajectory construction (root to leaf) cannot use virtual sub-states.
Only backward completion (leaf to root) can.
\end{theorem}

\begin{proof}
Forward construction at step $j$ produces $p_{j+1}$ as a refinement of $p_j$;
the system must commit to $p_{j+1} \in P_{j+1}$ as a physical block,
hence $\iota_{j+1}(p_{j+1}) \in [0,1]^3$.
Virtual sub-coordinates $(s_1,s_2,s_3) \notin [0,1]^3$ do not correspond
to any partition block and cannot serve as $p_{j+1}$ in a forward step.

Backward completion, by contrast, computes intermediate decompositions
\emph{analytically} within the receiver's internal calculus. The receiver's
evaluation map permits any decomposition satisfying mean-recovery
(Definition~\ref{def:sub_coord}); the partition requires only that
the global coordinate at each level lies in $[0,1]^3$.
The virtual sub-coordinates exist in the receiver's internal computation,
not in the physical partition.
\end{proof}

\subsection{Collapse without virtual states}

\begin{theorem}[Collapse Without Virtual States]
\label{thm:collapse}
Backward trajectory completion restricted entirely to physically realizable
sub-coordinates (all intermediate states in $[0,1]^3$) has worst-case
complexity $\Theta(N)$, identical to naive forward enumeration.
\end{theorem}

\begin{proof}
A physical-only backward step at depth $j$ requires identifying $p_{j-1}$
from $p_j$ without the geometric shortcut provided by the Fisher metric.
The algorithm must enumerate the children of candidate parent blocks to
verify which contains $p_j$.

At level $j$, there are $3^j$ candidate parent blocks; each has 3 children.
Without the embedded coordinates, verifying parenthood requires
$\mathcal{O}(3^j)$ comparisons per level. Summing over all levels:
$\sum_{j=0}^k 3^j = (3^{k+1}-1)/2 = \Theta(N)$.

More precisely: without geometric structure, the algorithm must construct
the hierarchy from scratch by queries, which requires $\Omega(N \log_2 3)$
bits of information (Theorem~\ref{thm:embedding_necessity}) and cannot be
done in $o(N)$ queries.
\end{proof}

\begin{corollary}[Virtual States are Necessary for Efficiency]
\label{cor:virtual_necessary}
Virtual sub-states are not optional features of backward trajectory completion.
Restricting to physical intermediate states increases the computational
cost from $\mathcal{O}(\log N)$ to $\Theta(N)$ — for $N = T(60,3) \approx
10^{59}$ LHC event categories, this is the difference between microseconds
and cosmological timescales.
\end{corollary}

% =========================================================
\section{Virtual Sub-States}
\label{sec:virtual}
% =========================================================

\subsection{Definition and existence}

\begin{definition}[Sub-Coordinate Decomposition]
\label{def:sub_coord}
A \emph{sub-coordinate decomposition} of a global coordinate $s \in [0,1]$
is an ordered triple $(s_1, s_2, s_3) \in \R^3$ satisfying the
\emph{mean-recovery constraint}:
\[
\bar s = \frac{s_1 + s_2 + s_3}{3} = s.
\]
The decomposition is \emph{physical} if $(s_1, s_2, s_3) \in [0,1]^3$.
It is \emph{virtual} if at least one $s_i \notin [0,1]$.
\end{definition}

\begin{theorem}[Virtual Sub-State Existence]
\label{thm:virtual_existence}
For every $s \in (0,1)$ and every $\varepsilon > 0$, the set of virtual
sub-coordinate decompositions with mean $s$ and
$\max_i |s_i| \leq M$ (for any $M \geq \varepsilon$) has positive
Lebesgue measure in the constraint plane
$\{(s_1,s_2,s_3) : (s_1+s_2+s_3)/3 = s\}$.
Moreover, the measure of virtual decompositions approaches 1 as $M \to \infty$.
\end{theorem}

\begin{proof}
The constraint plane is two-dimensional in $\R^3$. The physical region
(intersection with $[0,1]^3$) is a convex polygon of bounded area.
The virtual region (outside $[0,1]^3$) grows as $M^2$ with $M$.
The fraction of the constraint plane that is physical tends to 0 as $M \to \infty$,
so the virtual fraction tends to 1.
For any fixed $M \geq \varepsilon$, the specific set
$\{(s+\varepsilon, s+\varepsilon, s-2\varepsilon) : \varepsilon \in (0, M-s)\}$
contains virtual decompositions (first two components exceed 1 for
$\varepsilon > 1-s$) of positive measure.
\end{proof}

\begin{theorem}[Local-Global Decoupling for Trajectories]
\label{thm:local_global}
Let $\mathcal{T} = (p_0, \ldots, p_k)$ be a trajectory with endpoint
$\iota_k(p_k) = (\Sk, \St, \Se) \in [0,1]^3$.
For any target sequence of sub-coordinate decompositions
$((s_1^{(j)}, s_2^{(j)}, s_3^{(j)}))_{j=0}^k$
with mean-recovery $\bar s^{(j)} = \iota_j(p_j)_c$ (one component $c$
per level), there exists a receiver-internal computation realising
the trajectory with these sub-coordinates at each step.
\end{theorem}

\begin{proof}
Apply the Unconstrained Subtask theorem \citep{companion:recursion}
per step: the global S-value at step $j$ is the embedded coordinate;
the local sub-coordinate decomposition is a free choice subject to
mean-recovery; the cancellation pattern $(s_i - s_i)$ preserves the
global value. Construct level by level.
\end{proof}

\begin{corollary}[Miracle Trajectory]
\label{cor:miracle}
For any trajectory $\mathcal{T}$, there exists a receiver-internal
computation realising $\mathcal{T}$ in which \emph{every} intermediate
sub-coordinate is virtual (outside $[0,1]$).
\end{corollary}

\begin{proof}
At each step, choose sub-coordinates $(M+\varepsilon, M+\varepsilon,
3s - 2M - 2\varepsilon)$ with $M$ large. The first two components
exceed 1; the third is highly negative. The mean is $s$.
\end{proof}

\begin{remark}[Physical interpretation]
The Miracle Corollary has a striking physical interpretation.
A trigger that classifies an event as a $Z \to \mu^+\mu^-$ decay
can, internally, pass through sub-computations that are ``completely wrong''
(S-value = 100 for each individual step) while ultimately arriving
at the correct classification. This is not a failure of the algorithm —
it is the mechanism that makes efficient classification possible.
The virtual sub-states are not errors; they are the essential navigational
tools.
\end{remark}

\subsection{The miracle principle}

\begin{theorem}[Miracle Principle]
\label{thm:miracle}
For any target global S-value $\mathcal{S}^* \in [\Sfl(\receiver), 100]$
and any assignment of local S-values $(\sigma_1, \ldots, \sigma_n) \in
[\Sfl, 100]^n$, there exists a computation realising $\mathcal{S}^*$ globally
with $\mathcal{S}(\eta_i) = \sigma_i$ locally for each subtask $\eta_i$.
\end{theorem}

\begin{proof}
Construct the global computation as $\xi = \xi_0 + \sum_i (\eta_i - \eta_i)$
where $\xi_0$ has global S-value $\mathcal{S}^*$ and each $\eta_i$ is chosen
with local value $\sigma_i$. The cancellations $\eta_i - \eta_i = 0$
contribute nothing to the evaluation; the local S-values are by construction
the chosen targets.
\end{proof}

\begin{corollary}[Miraculous Subtasks]
Taking $\sigma_i = 100$ for all $i$ gives a globally correct computation
in which every subtask is maximally wrong locally. These are \emph{miraculous
subtasks}: individually wrong, collectively right.
\end{corollary}

% =========================================================
\section{Path Opacity and Its Trigger Implications}
\label{sec:opacity}
% =========================================================

\begin{theorem}[Path Opacity]
\label{thm:path_opacity}
Let $\mathcal{T}_1$ and $\mathcal{T}_2$ be two backward trajectories
sharing the same endpoint $p_k \in P_k$ but with distinct sequences of
virtual sub-coordinate decompositions at intermediate steps.
Given only the endpoint $p_k$ and the metric $g$, no metric invariant
computed at $p_k$ distinguishes $\mathcal{T}_1$ from $\mathcal{T}_2$.
\end{theorem}

\begin{proof}
By Theorem~\ref{thm:local_global}, intermediate sub-coordinate decompositions
are unconstrained subject to mean-recovery; any two trajectories with the
same endpoint but distinct intermediates both satisfy the trajectory
definition.

Every metric invariant $I(\mathcal{T}, g)$ that depends on the endpoint alone —
e.g., $\Svec(p_k)$, the geodesic distance $d(p_0, p_k)$, the Riemannian
volume of the constraint plane through $p_k$ — is by definition equal for
both trajectories.

Invariants depending on intermediate states are not accessible from the
endpoint; they cannot be computed from $p_k$ without additional information
about the path taken.
\end{proof}

\begin{corollary}[S-matrix Analog]
Path opacity is the categorical analog of S-matrix theory: only the initial
and final states (root and endpoint) are observable. The virtual
intermediate states are unobservable. This is not an accident: it is the
structural reason that quantum field theory's S-matrix can be computed from
Feynman diagrams that involve unphysical (off-shell, virtual) intermediate
states without contradiction.
\end{corollary}

\begin{remark}[Trigger implication]
Path opacity means that the trigger's internal computation — which virtual
sub-state decompositions it uses, which intermediate evaluations it performs —
is invisible to any external observer and irrelevant to the classification
result. The trigger may use any internally convenient virtual computations
as long as it arrives at the correct endpoint classification. This is the
formal justification for why the HLT can use approximations, look-up tables,
and cached partial results: the internal path does not affect the final
classification.
\end{remark}

\begin{theorem}[Non-Demolition Classification]
\label{thm:qnd_completion}
The backward completion algorithm makes exactly one observable: the root
category $p_0$. The intermediate virtual sub-states are computed
receiver-internally and never become physical observables.
The classification is therefore non-demolition: it extracts category
information from the endpoint without creating any observable disturbance
to the intermediate computation.
\end{theorem}

\begin{proof}
The backward algorithm (Section~\ref{sec:backward}) reads the trit string
of $\iota_k(p_k)$ and outputs $p_0$ (the root category). The virtual
sub-coordinates, which exist only in the receiver's internal S-space
computation, do not appear in any physical output. The classification
maps $p_k \mapsto p_0$ deterministically and completely; no disturbance
to $p_k$ or to the physical state at the endpoint is required or produced.
\end{proof}

% =========================================================
\section{The LHC Trigger as a Categorical Navigator}
\label{sec:trigger}
% =========================================================

\subsection{Event as endpoint, physics as root}

The LHC trigger classification problem maps exactly onto the trajectory
completion problem:

\begin{center}
\begin{tabular}{ll}
\toprule
\textbf{Trajectory completion} & \textbf{LHC trigger} \\
\midrule
Leaf $p_k$ (endpoint) & Detector hit pattern \\
Root $p_0$ & Originating physics process \\
Hierarchy depth $k$ & Number of detector subsystem levels \\
Backward navigation & Trigger classification \\
Virtual sub-states & Internal trigger computations \\
Path opacity & HLT internal computations invisible \\
$\mathcal{O}(\log_3 N)$ & Real-time classification in $<2.5\,\mu$s \\
\bottomrule
\end{tabular}
\end{center}

\begin{definition}[LHC Event Hierarchy]
\label{def:lhc_hierarchy}
The LHC event classification hierarchy is a ternary refinement hierarchy
of depth $k = 3$ (current ATLAS) to $k = 4$ (HL-LHC) with levels:
\begin{enumerate}[nosep,label=(\arabic*)]
\item Level 0 (root): broad physics category (QCD background, electroweak, Higgs).
\item Level 1: specific process type (dijet, $W$/$Z$ production, $t\bar t$, Higgs).
\item Level 2: decay channel ($W \to \ell\nu$, $Z \to \ell\ell$, $H \to \gamma\gamma$).
\item Level 3 (leaf): specific detector hit pattern (trigger condition fired).
\end{enumerate}
Each level trifurcates: three broad categories at level 0, nine at level 1, etc.
\end{definition}

\subsection{Cell registry as the S-entropy embedding}

\begin{theorem}[Physical Origin of the Trigger Cell Registry]
\label{thm:cell_registry}
The trigger's cell registry $\Gamma = \{C_i\}$ is the physical realisation
of the S-entropy embedding $\iota_k$: each cell $C_i$ corresponds to a
cell in the $3^k$-fold partition of $[0,1]^3$, and the partition coordinates
$(n,\ell,m,s)$ of the detected particle determine which cell $C_i$ it
falls into.
\end{theorem}

\begin{proof}
From \citet{companion:math} and \citet{companion:planck}, the partition
coordinates $(n,\ell,m,s)$ label the cells of the ternary refinement
hierarchy at depth $k = 4$ (for the four subsystems tracker/$n$, ECAL/$\ell$,
HCAL/$m$, muon/$s$). The S-entropy embedding maps each combination of
partition coordinates to a unique point in $[0,1]^3$ via the base-3
expansion of Definition~\ref{def:s_embed}.

The trigger's timing cell $C_i = \{|\Delta P| < \delta_i\}$ constrains
the $\St$ coordinate of the event (temporal entropy — how close the signals
are to the reference bunch crossing). Combined with the ECAL shower shape
(constraining $\Se$) and track momentum (constraining $\Sk$), the three
coordinates together uniquely address a leaf in the S-entropy cube.
\end{proof}

\subsection{The di-muon trigger as a worked example}

The di-muon trigger ($Z \to \mu^+\mu^-$ selection) is the simplest
physics-motivated trigger and the natural starting point for illustration.

\begin{example}[Di-Muon Backward Completion]
\label{ex:dimuon}
A collision event produces two muons of opposite sign with $p_T > 20\,\mathrm{GeV}$.
The backward completion proceeds:

\emph{Level 3 (endpoint)}: The muon system fires for two tracks in opposite
hemispheres. In the S-entropy embedding:
$\Sk \approx 0.7$ (high momentum tracks, near-maximum energy shell),
$\St \approx 0.5$ (signals arrive on-time, at the median temporal entropy),
$\Se \approx 0.3$ (low evolution entropy, clean two-track topology).

\emph{Level 2 (backward step 1)}: Parent cell has $\Se \approx 0.33$
(slightly higher: includes the post-decay kinematics). This corresponds to
the $Z \to \ell^+\ell^-$ decay channel.

\emph{Level 1 (backward step 2)}: Parent cell has $\Se \approx 0.5$
(intermediate evolution: the $Z$ boson produced in the collision).
This corresponds to Drell-Yan $Z$ production.

\emph{Level 0 (root)}: $\Se \approx 0.8$ (high evolution entropy: the
original quark-antiquark annihilation). This is the electroweak production category.

The three backward steps ($k = 3$) recover the complete physics process
from the detector hit pattern in exactly 3 operations.

The virtual sub-states are the intermediate decompositions:
at Level 2 $\to$ Level 1, the computation internally uses
$(s_1, s_2, s_3) = (0.6, 0.4, -0.1)$ for the $\Se$ coordinate
(the $-0.1$ is virtual). The mean $(-0.1 + 0.6 + 0.4)/3 \approx 0.3$
recovers the physical value. The virtual component is invisible to any
external measurement.
\end{example}

\subsection{Why selection rules are structural dead zones}

From the partition coordinate structure \citep{companion:math}, selection
rules $\Delta\ell = \pm 1$, $\Delta m \in \{0,\pm 1\}$ constrain which
transitions between partition states are physically realizable. In the
backward completion framework, this has a direct consequence:

\begin{theorem}[Selection Rules as Type Errors]
\label{thm:selection_rules}
Trigger backgrounds that require selection-rule-violating intermediate
transitions between partition states are \emph{type errors} in the vaHera
fragment representing the event. They are rejected by the pre-dispatch
validation engine at static time before any backward completion is attempted,
and they have no pre-image under the backward parent map from any physically
realizable endpoint.
\end{theorem}

\begin{proof}
In the refinement-type system over the vaHera AST
\citep{companion:types}, every detector subsystem is a refined operation
whose output type carries the partition coordinate it measures, and whose
output refinement predicate is the conjunction of all registered selection
rules. A vaHera Compose chain whose implied transition violates any selection
rule has an uninhabited refinement type (the Forbidden-Transition Theorem of
\citealt{companion:types}) and is rejected by $\mathrm{typecheck}$ at static
time before execution.

In the ternary hierarchy, a transition $\Delta\ell = \pm 2$ (for example)
would require a parent cell at level $j-1$ that is two angular levels removed
from the child cell at level $j$. By the partition capacity formula
$C(n) = 2n^2$ and the $SO(3)$ structure of the partition coordinates
\citep{companion:math}, no two adjacent levels in the refinement hierarchy
are connected by $|\Delta\ell| \geq 2$ transitions. Therefore, no backward
step from a physical leaf $p_k$ can traverse such a transition: it would
require a non-existent parent block.

The two arguments are equivalent and reinforce each other: the
partition-hierarchy argument shows the background has no pre-image under
the backward parent map; the type-theoretic argument shows that any vaHera
expression attempting to construct such a transition is syntactically
ill-typed. Forbidden backgrounds are not merely rare --- they cannot be written.
\end{proof}

\begin{remark}[Type error vs.\ probabilistic suppression]
The difference between Theorem~\ref{thm:selection_rules} and conventional
background suppression is categorical: conventional triggers suppress
backgrounds by imposing energy thresholds (probabilistic); the partition
framework excludes backgrounds by structural non-existence (categorical).
A background requiring $\Delta\ell = 2$ is not ``rare'' — it is absent from
the backward completion tree entirely.
\end{remark}

% =========================================================
\section{The Aperture Identification Theorem}
\label{sec:aperture}
% =========================================================

\subsection{Five names for one object}

The following five descriptions appear in different branches of physics,
mathematics, and computer science. We prove they refer to the same
mathematical object.

\begin{theorem}[Aperture Identification]
\label{thm:aperture_id}
The following five descriptions name the same mathematical object,
which we call the \emph{categorical aperture}:

\begin{enumerate}[nosep,label=(\arabic*)]
\item \textbf{Geometric aperture}: the convex hull of admissible sub-coordinate
decompositions of a fixed global coordinate, intersected with the virtual
region $\R^3 \setminus [0,1]^3$.

\item \textbf{Virtual sub-state}: a non-physical intermediate decomposition
$(s_1, s_2, s_3) \notin [0,1]^3$ with mean $s \in [0,1]$, used in
backward trajectory completion (Section~\ref{sec:virtual}).

\item \textbf{Miracle subtask}: a subtask $\eta$ in a computation with
local S-value $\mathcal{S}(\eta) = 100$ (maximally misaligned) but
contributing to a globally aligned composition
(Theorem~\ref{thm:miracle}).

\item \textbf{Information catalyst}: an expression atom whose application
strictly decreases the receiver's S-distance to truth,
$\mathcal{S}(\xi \diamond \gamma, x^*) < \mathcal{S}(\xi, x^*)$
\citep{companion:recursion}.

\item \textbf{Corrected Maxwell demon}: a sorting element that selects
which partition cells are accessible based on their categorical addresses,
operating without thermodynamic work (from Axiom~\ref{ax:bounded},
the categorical observable commutes with the physical Hamiltonian).
\end{enumerate}

All five share four defining properties:
enable a transition; not consumed in enabling it; transfer zero information
about themselves to the transition outcome; are necessary for the
$\mathcal{O}(\log N)$ speedup.
\end{theorem}

\begin{proof}
We exhibit explicit identifications between each pair.

\emph{(1) $\leftrightarrow$ (2)}: The convex hull of admissible decompositions
of a global coordinate $s$ intersected with $\R^3 \setminus [0,1]^3$ is
precisely the set of virtual sub-states with mean $s$. This is immediate
from Definition~\ref{def:sub_coord}.

\emph{(2) $\leftrightarrow$ (3)}: A virtual decomposition
$(s_1,s_2,s_3) \notin [0,1]^3$ has at least one component $s_i$ that,
as an independent sub-coordinate, would have $\mathcal{S}(s_i, x^*) = 100$
(it lies outside the physical state space, maximally misaligned). This is
the miracle subtask: individually $\mathcal{S} = 100$, collectively the
mean recovers the correct value.

\emph{(3) $\leftrightarrow$ (4)}: A catalyst $\gamma$ is an expression
whose application reduces the global S-value. The Miracle Principle
(Theorem~\ref{thm:miracle}) gives catalyst-like constructions; conversely,
any catalyst can be realised as a virtual sub-state decomposition
(Theorem~\ref{thm:local_global}).

\emph{(4) $\leftrightarrow$ (5)}: An information catalyst that operates on
categorical addresses rather than physical quantities is precisely a
sorting element that uses geometric (positional) information about
categorical cells to direct transitions. By the non-demolition theorem
(Theorem~\ref{thm:qnd_completion}), this sorting requires zero thermodynamic
work (it operates on the categorical degree of freedom, not the physical
degree of freedom). This is the corrected Maxwell demon: not a paradox,
but a physical sorting element that exploits the commutativity
$[\hat{O}_{\mathrm{cat}}, \hat{O}_{\mathrm{phys}}] = 0$
from \citet{companion:planck}.

\emph{Shared properties}: All five enable transitions by providing geometric
shortcuts through state space. None is consumed: the virtual sub-state
exists as long as its mean-recovery constraint is satisfied; the catalyst
can be applied multiple times. None transfers information about itself:
path opacity (Theorem~\ref{thm:path_opacity}) shows the virtual path is
invisible. All are necessary: Theorem~\ref{thm:collapse} shows that without
virtual states, complexity collapses to $\Theta(N)$.
\end{proof}

\begin{corollary}[The LHC Trigger is a Categorical Aperture]
\label{cor:trigger_aperture}
The LHC trigger implements the categorical aperture. Its timing cells define
which S-entropy regions are accessible (geometric aperture). Its internal
HLT computations use virtual intermediate states (virtual sub-states). Its
trigger conditions are individually approximate (miracle subtasks: each
condition has non-zero false positive rate). The trigger collectively achieves
precise event classification (globally aligned composition). It operates
without disturbing the classified event (non-demolition). And it enables
real-time classification that would be impossible by forward enumeration
(necessary for the $\mathcal{O}(\log N)$ speedup).
\end{corollary}

% =========================================================
\section{Discussion}
\label{sec:discussion}
% =========================================================

\subsection{Complexity summary}

The main complexity results are:
\begin{itemize}[nosep]
\item Backward completion with virtual sub-states: $\mathcal{O}(\log_3 N)$
(Theorem~\ref{thm:backward_complexity}).
\item Backward completion restricted to physical states: $\Theta(N)$
(Theorem~\ref{thm:collapse}).
\item Forward enumeration (conventional trigger approach): $\Theta(N)$.
\end{itemize}
For $N = T(60,3) \approx 10^{59}$ LHC event categories, the difference
between $\log_3(10^{59}) = 98$ steps and $10^{59}$ steps is the difference
between real-time trigger decisions and physically impossible computation.

\subsection{The S-entropy floor as trigger precision floor}

The Floor Positivity theorem (Theorem~\ref{thm:floor}) has a direct trigger
interpretation. No bounded trigger system can achieve $\mathcal{S}(\text{trigger}(x), x^*) = 0$
for all events: the trigger's finite knowledge framework $K$ cannot represent
every possible event topology. The floor $\Sfl(\receiver) > 0$ quantifies
the irreducible misclassification probability of any trigger of given complexity.

For the ATLAS L1 trigger with $\sim 10^3$ trigger paths out of $\sim 10^{59}$
possible topologies, the composition efficiency $\etaC \approx 0.08$ from
\citet{companion:planck} quantifies how far the trigger is from the theoretical
floor. Increasing $\etaC$ toward 1 (filling the $n_P = 60$ depth) would
reduce the floor and decrease the irreducible misclassification probability.

\subsection{Relation to Kalman filtering}

The HLT track reconstruction uses the Kalman filter, which is a forward-backward
algorithm: forward pass accumulates state estimates; backward pass (smoothing)
corrects them using future measurements. The backward pass is precisely
backward trajectory completion in the particle trajectory state space.
The virtual intermediate states in the backward completion are exactly the
extrapolated state estimates that lie outside the detector acceptance in
intermediate layers — ``virtual'' in the sense that no detector exists at
that position, yet essential for the smoothed estimate.

This observation suggests that Kalman filtering and S-entropy backward
completion are the same algorithm in different coordinate systems, related
by the Triple Equivalence (Theorem~\ref{thm:triple_equiv}).

\subsection{Open questions}

\textbf{Path multiplicity}: For a fixed endpoint $p_k$, how many distinct
backward trajectories exist in the ternary hierarchy? By the Path Opacity
theorem, they are all equivalent from the observer's perspective, but their
total count bounds the trigger's internal computational freedom.

\textbf{Optimal virtual state selection}: Among all virtual sub-state
decompositions achieving mean-recovery, which minimise the Fisher information
along the backward path? This would give an optimal backward navigation
algorithm rather than the generic one of Algorithm~1.

\textbf{Cascade catalysis and trigger iterations}: The multiplicative catalyst
algebra (proved in \citet{companion:cascade}) says catalyst power combines
as $\catpow(\gamma_1 \diamond \gamma_2) = 1 - (1-\catpow_1)(1-\catpow_2)$.
A trigger menu of $K$ trigger paths is a catalyst cascade; the total
classification power is $1 - \prod_i (1 - \catpow_i)$. Saturating the
cascade ($\catpow \to 1$) requires $\sum_i \catpow_i = \infty$, which
occurs when the trigger menu's composition depth reaches $\nP$.

\section*{Acknowledgements}
This work was supported by the AIMe Registry for Artificial Intelligence.

\bibliographystyle{plainnat}
\bibliography{references}

\end{document}
